\section{Conclusion}
\label{sec:conclusion}

The PT zeta programme delivers a canonical spectral model of the
non-trivial zeros of the Riemann zeta function from a single internal
structure: the persistence dynamics with its axiom $\sper = 1/2$.
The final balance can be summarised in five points.

\begin{enumerate}[leftmargin=*]
\item \textbf{The canonical operatorial object is unique.} The
regularised trace
\[
\Treg(s) \;=\; \Tr_{\mathrm{reg}}\!\big[(s - i\HPTBK)^{-1}\,\DRop(s)\big]
\;=\; \sum_{p \in \Primes_\gamma} \kappap(s)\,(\log p)\,\Big[\frac{1}{p^s - 1} + \Ap\,p^{-s}\Big]
\]
is constructed from the amplitudes $\kappap$ of the
$\qplus/\qminus$ bifurcation and the PT Berry--Keating dilation
operator. No free parameter. Exact analytic identity in
$\Reop(s) > 1$ to $10^{-26}$, distributional poles on
$\Reop(s) = 1/2$ at the zeros of
$\Rres(s) \equiv \zeta(s) / (\zetaplus\,\zetaminus)$.

\item \textbf{The critical line $\Reop(s) = 1/2$ admits four
equivalent readings}: symplectic Planck cell
$u_\star\,p_\star = 2\pi$, multiplicative Haar Jacobian, unitary
transformation $e^{su/2}$, antiperiodic boundary condition forced
by $T_3$. The three roles of $\sper = 1/2$ ($\Tone$ arithmetic,
$\Tfive$ geometric, Haar spectral) are operatorially identified.

\item \textbf{The cohomological residue $1/8$ is triply determined.}
\[
\dfrac{1}{8} \;=\; \dfrac{\sper^{2}}{2} \;=\; \lambdaH
\;=\; \dfrac{c_1}{N_{\mathrm{corners}}} \;=\; \dfrac{\sper}{4}.
\]
The coherence $\sper^{2}/2 = \sper/4$ admits $\sper = 0$ or
$\sper = 1/2$; $\Tone$ excludes $\sper = 0$, which overdetermines
$\sper = 1/2$: this is a structural signature of $\Tone$ by the
HP-PT zone, not a substitute for $\Tone$.

\item \textbf{Structural Higgs $\leftrightarrow \zeta$ duality.}
The canonical bifurcator $B$ (spinor projector $\qplus/\qminus$ of
the channel $p = 2$ at the fixed point $\mustar = 15$) produces
simultaneously $m_H = v \cdot \sper$ and
$\Delta_N = \lambdaH = \sper^{2}/2$ via the chain
$\Tone$ + $\Ttwo$ + Haar + Weyl. No other PT coupling equals
$1/8$: pure numerical coincidence is rejected.

\item \textbf{Falsifiability discipline.} The Dirichlet prediction
$\Delta_N(L \bmod q) = 1/(8q)$, structurally induced by K4, was
tested on LMFDB zeros of four primitive quadratic characters and
falsified ($\chi^{2}(\mathrm{PT}) = 7.49$ against
$\chi^{2}(H_0) = 0.023$ on 4~d.o.f.). The residue $1/8$ is an
asymptotic identity between smooth forms, not a measurable shift in
the zero count. The PT mechanism described here is specific to
$\zeta$.
\end{enumerate}

\paragraph{Canonical position.}
We have \textbf{understood} the zero structure of $\zeta$; we have
\textbf{not proved} their strict localisation on the critical line.
The ultimate lock ($G_{\mathrm{HP3.a}}$: strict analytic continuation
of $\Rres$) is equivalent to RH in the classical sense and does not
seem any easier than the latter.

The main enrichment over Berry--Keating--Connes is the intrinsic
reading of the coincidences: the $2\pi$ is the canonical commutator,
the $1/8$ is the scalar Higgs self-coupling of PT, and the spin
$\sper = 1/2$ of the critical line is overdetermined by the
cohomological coherence of the construction. \textbf{No ingredient
is ad hoc.}
